\scanpage{205}
Ainsi, $\pi_1(X)$ n'opère évidemment pas sur $\xi_{x,\mathbf Z}$, \unsure{ni donc} sur $\xi_{x,\mathbf R}$, mais l'élément $f_\infty$ de $\pi_1(X)$ (frobenius infini) opère sur $\xi_{x,\mathbf Z}$ et à fortiori sur $\xi_{x,\mathbf R}$ ; on le prolonge \uline{linéairement} en $\tau=f_\infty$ ....

\uline{N.B.} La \uline{conjecture de Hodge} interprète la filtration \uline{intrinsèque} sur \unsure{la cat. réalisée banalement} $\xi_{x,\mathbf Z}\otimes_{\mathbf Z}\mathbf Q$ \inserted{$=\xi_{x,\mathbf Q}$} par la filtration arithmétique ($=$ motivique, par le niveau), en termes de la position par. à la graduation [i.e. en termes de la filtration de Hodge -- non motivique -- et du $f_\infty$]. Elle s'exprime en disant que

\deleted{$F$ de niveau $\leq\unsure{\rho}$ $\Longleftrightarrow$}

\uline{La partie de niveau} $\leq\unsure{\rho}$ \uline{dans} $\xi_{x,\mathbf Q}(F)$ \uline{n'est autre que}
\[
\xi_{x,\mathbf Q}(F)\cap\xi_{x,\infty}(F)^{\unsure{\rho}}\cap f_\infty\xi_{x,\infty}(F)^{(\unsure{\rho})}.
\]

En particulier, le dernier sous-espace de $\xi_{x,\mathbf Q}(F)$ est \uline{motivique} (provient d'un sous-motif $F^{(\unsure{\rho})}$ de $F$), alors que $\xi_{x,\infty}(F)^{(\unsure{\rho})}$ et $f_\infty\xi_{x,\infty}(F)^{(\unsure{\rho})}$ n'est pas motiviques .... \deleted{Si $\unsure{\rho}=0$, \unsure{cela donne} l'intersection des \unsure{réseaux} ``\unsure{normaux}''}

\mspage{204}
